\documentclass[../main.tex]{subfiles}
\begin{document}
\section{Detalhamento de Serviços e Arquitetura Operacional}

A \textbf{CSIS} abdica da estrutura engessada de um produto de software tradicional, convergindo para um "Ecossistema de Serviços e Educação". Esse modelo opera na interseção do desenvolvimento contínuo de talentos (B2C) e na prestação de perícia técnica avançada (B2B), criando um ciclo autossustentável: a formação prática rentável subsidia a força de trabalho investigativa de altíssimo padrão com custos marginais declinantes.

\subsection{O Catálogo Forense e Corporativo (B2B)}
No escopo de mercado voltado para empresas e escritórios de advocacia, a CSIS comercializa inteligência técnica digital tangível e homologada. O portfólio de serviços inclui:
\begin{itemize}
    \item \textbf{Laudos de Perícia Digital e Computação Forense:} Extrações, processamento e análises (\textit{Forensics}) abrangendo dispositivos móveis e infraestruturas computacionais corporativas. O serviço foca em encontrar artefatos digitais fidedignos, indispensáveis para instrução de litígios.
    \item \textbf{Testes de Invasão (\textit{Pentest}):} Execução tática de simulações (\textit{Red Team}) para identificação cirúrgica e exploração controlada de vulnerabilidades em sistemas empresariais, atuando massivamente na prevenção de resgates maliciosos (\textit{ransomware}) e neutralizando vazamentos de rede.
    \item \textbf{Consultoria de Mitigação de Riscos:} Pareceres táticos avançados para times de Defesa (\textit{Blue Team}) orientando pequenas e médias empresas a blindarem processos alinhados à legislação atual.
\end{itemize}

Qualquer dossiê ou laudo destinado à ponta B2B passa, impreterivelmente, através da rígida câmara final de acompanhamento e revisão de Excelência (\textit{QA}) do Perito Sênior, dotando as peças da indispensável credibilidade jurídica em tribunais ou diretorias corporativas.

\subsection{O Produto Acadêmico e Repositório (B2C)}
Para sustentar o exército da análise operacional com excelência sem originar enormes passivos contratuais e previdenciários de equipe estática, atua-se na distribuição de um infoproduto imersivo. O talento (estudante B2C) não adquire apenas teoria engessada, ele acessa o "Ecossistema Tático":
\begin{itemize}
    \item \textbf{Acesso Contínuo e Direcionado ao Currículo:} Transmissão de conhecimento categorizada estrategicamente em fluxos modulares progressivos de complexidade técnica (Fase \textit{Childrens}, Fase \textit{Basic} e Fase \textit{Advanced}).
    \item \textbf{Módulo Especializado de \textit{Report Writing}:} Ponto-cego severo da maioria dos cursos de Segurança, essa etapa é fundamental na CSIS. Foca o pesquisador na montagem visual e escrita das provas forenses perante leigos técnicos (\textit{Proof of Concept} e relatórios compreensíveis à juízes/gestores).
    \item \textbf{O Círculo de Monetização (\textit{Bounties}):} O propulsor do ecossistema. Alunos diplomados nos desafios internos ganham chancela para capturar "Missões Remuneradas" que solucionam frações reais dos casos do núcleo B2B, permitindo rentabilidade estudantil simultânea à construção material empírica (\textit{Portfólio Valioso}).
\end{itemize}

\subsection{Arquitetura e Redução Drástica de Passivo Físico}
Operando com um contorno \textit{Bootstrapped-Lean}, as tecnologias estruturais da CSIS não abraçam o \textit{over-engineering} (manutenção de plataformas acadêmicas exclusivas), centralizando-se somente no pragmatismo nativo das grandes \textit{big techs} forenses:
\begin{itemize}
    \item \textbf{Aquisição e Mentoria:} Hospedagem robusta da grade de vídeo aulas suportando o \textit{streaming} via \textit{Hotmart} (ou análogas); aquisição orgânica de autoridade suportada pelo topo de funil no formato de tutorais vitrines do \textit{YouTube}.
    \item \textbf{Orquestramento de \textit{Bounties} e Colaboradores:} Integração técnica feita pela matriz relacional no Discord — gerenciando hierarquia humana dos estagiários — enquanto todo versionamento do ciclo de Missão para Laudo Técnico ocorre limpidamente em \textit{Issues} e \textit{Pull Requests} privados dentro do ecossistema do \textit{GitHub}.
    \item \textbf{Laboratórios Remotos Não-Hospedados:} A empresa delibera como política operacional inicial que as emulações e laboratórios virtuais rodem descentralizados na ponta do hardware do próprio talento acadêmico (alunos \textit{freelancers}), esmagando a necessidade contínua de hospedagens maciças e aquecimento de infraestrutura nas Nuves empresariais (\textit{Cloud Cost Mitigation}).
\end{itemize}

No organograma tático na figura \ref{fig:fluxo_infoproduto}, elucida-se a circularidade perfeita entre desenvolvimento da força braçal técnica B2C e a propulsão financeira comercial pericial da Boutique.

\begin{figure}[h]
    \centering
    % NOTA PARA O AUTOR: O latex quase sempre falha quando há PDFs com espaços e parênteses. 
    % Recomendo fortemente renomear o seu arquivo para 'Cyber_Infoproduto_Talent.pdf' na pasta /img. 
    \includegraphics[width=\textwidth]{img/Cyber_Infoproduto_Talent.pdf}
    \caption{A Roda Virtuosa "Talent Flywheel": Fluxo simbiótico entre o infoproduto educacional e a formação orgânica chancelada.}
    \label{fig:fluxo_infoproduto}
\end{figure}

\end{document}